% Standalone problem document, generated from the adjacent README.md.
% Compile with XeLaTeX. Mathematical content is maintained in Markdown.
\documentclass[11pt,a4paper]{article}
\usepackage[fontsize=10.5pt]{scrextend}
\usepackage[margin=22mm,headheight=15pt,headsep=9mm,footskip=11mm]{geometry}
\usepackage{fontspec}
\setmainfont{texgyrepagella}[Extension=.otf,UprightFont=*-regular,BoldFont=*-bold,ItalicFont=*-italic,BoldItalicFont=*-bolditalic]
\setsansfont{texgyreheros}[Extension=.otf,UprightFont=*-regular,BoldFont=*-bold,ItalicFont=*-italic,BoldItalicFont=*-bolditalic]
\setmonofont{texgyrecursor}[Extension=.otf,UprightFont=*-regular,BoldFont=*-bold,ItalicFont=*-italic,BoldItalicFont=*-bolditalic,Scale=MatchLowercase]
\usepackage{amsmath,amssymb}
\usepackage{unicode-math}
\setmathfont{texgyrepagella-math.otf}
\usepackage{xcolor,microtype,enumitem,needspace,fancyhdr}
\usepackage{bookmark}
\definecolor{ink}{HTML}{17344B}
\definecolor{accent}{HTML}{197D80}
\definecolor{muted}{HTML}{596773}
\hypersetup{colorlinks=true,linkcolor=accent,urlcolor=accent,pdftitle={MF-22 - Polynomial
conditioning of the cubic C1 spline Schrödinger Toeplitz
family},pdfauthor={Open Problems in Numerical Linear Algebra}}
\urlstyle{same}
\setlength{\parindent}{0pt}
\setlength{\parskip}{5pt plus 1pt minus 1pt}
\linespread{1.04}
\setlength{\emergencystretch}{3em}
\widowpenalty=10000
\clubpenalty=10000
\displaywidowpenalty=10000
\allowdisplaybreaks[1]
\setlist{nosep,leftmargin=1.5em}
\providecommand{\tightlist}{\setlength{\itemsep}{0pt}\setlength{\parskip}{0pt}}
\providecommand{\pandocbounded}[1]{#1}
\pagestyle{fancy}
\fancyhf{}
\fancyhead[L]{\sffamily\footnotesize\color{muted}OPEN PROBLEMS IN NUMERICAL LINEAR ALGEBRA}
\fancyhead[R]{\sffamily\bfseries\footnotesize\color{accent}MF-22}
\fancyfoot[L]{\parbox[b]{0.9\textwidth}{\sffamily\fontsize{8}{10}\selectfont\color{muted}Editorial ratings; a bounded search cannot certify that a problem remains unsolved.\\Literature check: 2026-09-22}}
\fancyfoot[R]{\sffamily\footnotesize\color{muted}\thepage}
\renewcommand{\headrulewidth}{0.3pt}
\renewcommand{\headrule}{\hbox to\headwidth{\color{accent}\leaders\hrule height \headrulewidth\hfill}}
\makeatletter
\renewcommand{\section}{\Needspace{5\baselineskip}\@startsection{section}{1}{0pt}{12pt plus 2pt minus 2pt}{4pt}{\sffamily\bfseries\large\color{ink}}}
\renewcommand{\subsection}{\Needspace{4\baselineskip}\@startsection{subsection}{2}{0pt}{9pt plus 1pt minus 1pt}{3pt}{\sffamily\bfseries\normalsize\color{ink}}}
\makeatother
\setcounter{secnumdepth}{0}
\begin{document}
{\sffamily\small\color{muted}Matrix functions and stability\par}
\vspace{3pt}
{\raggedright\hyphenpenalty=10000\exhyphenpenalty=10000\sffamily\bfseries\fontsize{21}{24}\selectfont\color{ink}Polynomial
conditioning of the cubic C1 spline Schrödinger Toeplitz family\par}
\vspace{7pt}
{\sffamily\small\textbf{Difficulty:} hard\quad\textbf{Importance:} interesting
to specialist\par}
{\sffamily\small\bfseries\color{accent}Status: Lean verified\par}
\vspace{4pt}
{\color{accent}\hrule height 0.8pt}
\vspace{6pt}
\textbf{Topic:} Block Toeplitz systems and discretization stability

\textbf{Rating rationale:} Historical ratings retained. The fixed
explicit band structure makes this a hard asymptotic question, but
available symbol criteria leave the relevant regime undecided; its
immediate importance is to specialists in structured discretization
stability.

\subsection{Lean verification --- 22 September
2026}\label{lean-verification-22-september-2026}

For every fixed real \(\rho>0\), the exact uncorrected complex
\(2n\times2n\) Toeplitz family is nonsingular for all sufficiently large
integer sizes and satisfies \(\kappa_2(H_n)\le K_\rho n\). Constants and
the size threshold may depend on the parameter, but not on the size. The
exceptional parameter \(\rho^2=10\) is included. Norms are the genuine
induced Euclidean operator norms, and the inverse is the actual matrix
inverse; singular matrices have infinite extended condition number. This
proves the complete original polynomial-growth target with exponent one.
The proof also establishes a dimension-independent forward norm bound
and eventual uniform bounds on every inverse entry.

All four exports passed LeanCert kernel audits, two independent
nonauthor final reviews and
\href{https://github.com/sidneyholden1/OpenProblemsInNLA/actions/runs/35707732163}{actual
isolated Linux Comparator/default-kernel verification}, including
rejection and sandbox controls.
\href{https://github.com/sidneyholden1/OpenProblemsInNLA/blob/main/matrix-functions-and-stability/MF-22/lean/README.md}{Proof
and reviews} ·
\href{https://github.com/sidneyholden1/OpenProblemsInNLA/tree/5596784dc81e22ed625650bca71d4a1435896cef/matrix-functions-and-stability/MF-22/lean}{Immutable
proof revision} ·
\href{https://github.com/sidneyholden1/OpenProblemsInNLA/blob/main/docs/lean/verification-2026-09-22/MF-22/README.md}{Retained
evidence}.

Mathematics: George Stepaniants, Caltech, with source-disclosed AI
assistance. Bogoya, Böttcher, Ferrari, Grudsky and Serra-Capizzano
retain original family/question credit. Formalization: Sidney Holden
with OpenAI Codex assistance. No novelty, source-author endorsement or
external human peer review is claimed. The complete original statement,
earlier resolution and source record remain below unchanged.

\pagestyle{plain}

\subsection{Resolution: affirmative, 11 September
2026}\label{resolution-affirmative-11-september-2026}

\textbf{George Stepaniants, Department of Computing and Mathematical
Sciences, California Institute of Technology}, proves that for every
fixed \(\rho>0\) there are \(K_\rho>0\) and \(n_\rho\) such that

\[\kappa_2(H_n(\rho))\le K_\rho n\qquad(n\ge n_\rho).\]

Thus the original question is answered affirmatively with exponent
\(\alpha_\rho=1\). The result includes eventual invertibility, the
parameter \(\rho=\sqrt{10}\), and the exact uncorrected Toeplitz
boundary entries. The original ID, full statement, and historical
ratings are retained below.

The precise locator is the \textbf{Theorem in Section 1, proved in
Sections 2--4} of the
\href{https://github.com/ajt60gaibb/OpenProblemsInNLA/blob/main/matrix-functions-and-stability/MF-22/solution.md}{complete
manuscript}.
\href{https://github.com/ajt60gaibb/OpenProblemsInNLA/blob/main/matrix-functions-and-stability/MF-22/solution.pdf}{Proof
PDF} ·
\href{https://github.com/ajt60gaibb/OpenProblemsInNLA/blob/main/matrix-functions-and-stability/MF-22/solution.tex}{Standalone
TeX}. The proof derives a four-state recurrence, excludes all numerator
cancellation, and bounds its finite Green matrix to obtain a linear
inverse bound.

A separate
\href{https://github.com/ajt60gaibb/OpenProblemsInNLA/blob/main/references/stepaniants-mf22-2026-09-11/verification/MF-22-independent-review.md}{independent
Codex-agent mathematical review} returned \textbf{PASS} for the complete
target. The
\href{https://github.com/ajt60gaibb/OpenProblemsInNLA/blob/main/references/stepaniants-mf22-2026-09-11/README.md}{submission
record} documents substantial ChatGPT/Codex assistance, exact symbolic
checks, the public branch/fork audit, and verification limits. This is
independent automated-agent review, not external human peer review or
formal proof certification. The source authors retain credit for the
family, the root classification and the question.

\subsection{Statement}\label{statement}

Define real \(2\times2\) matrices by

\[\begin{aligned}
B_{-1}&=\frac3{40}\begin{pmatrix}-1&0\\-5&0\end{pmatrix},&
B_0&=\frac3{40}\begin{pmatrix}0&-5\\16&-5\end{pmatrix},\\
B_1&=\frac3{40}\begin{pmatrix}-5&16\\-5&0\end{pmatrix},&
B_2&=\frac3{40}\begin{pmatrix}0&-5\\0&-1\end{pmatrix},\\
C_{-1}&=\frac1{80}\begin{pmatrix}1&0\\7&0\end{pmatrix},&
C_0&=\frac1{80}\begin{pmatrix}24&7\\0&25\end{pmatrix},\\
C_1&=\frac1{80}\begin{pmatrix}-25&0\\-7&-24\end{pmatrix},&
C_2&=\frac1{80}\begin{pmatrix}0&-7\\0&-1\end{pmatrix}.
\end{aligned}\]

Set \(B_k=C_k=0\) for all other integer indices. For a real parameter
\(\rho>0\) and \(n\ge1\), form the \(2n\times2n\) complex block Toeplitz
matrix

\[H_n(\rho)=\bigl(iB_{j-k}-\rho C_{j-k}\bigr)_{j,k=0}^{n-1},
\qquad i^2=-1.\]

Use the spectral condition number \(\kappa_2(H)=\|H\|_2\|H^{-1}\|_2\)
for nonsingular \(H\) and \(\kappa_2(H)=+\infty\) for singular \(H\).

\textbf{Question (explicit asymptotic formulation of the source's
polynomial-growth question).} For every fixed \(\rho>0\), do there exist
constants \(K_\rho>0\), \(\alpha_\rho\ge0\), and \(n_\rho\in\mathbb N\),
independent of \(n\), such that

\[\kappa_2(H_n(\rho))\le K_\rho n^{\alpha_\rho}
\qquad\text{for every }n\ge n_\rho?\]

The parameter remains fixed as \(n\to\infty\); no uniform-in-\(\rho\)
bound or particular exponent is asserted. Eventual invertibility is part
of the question. This is the exact pure Toeplitz family denoted
\(\widetilde{\mathbf S}_n^{(3,1)}(\rho)\) in the source; \(n\) here
counts its \(2\times2\) blocks. It has no additional corner corrections.

\subsection{Numerical significance}\label{numerical-significance}

These matrices arise from cubic splines with \(C^1\) continuity in time
for a space-time Schrödinger discretization. The source proves a
determinant-root classification of the associated polynomial symbol but
expressly leaves polynomial growth of these condition numbers open. The
linear bound recorded above explains the numerical stability seen for
this intermediate spline regularity.

\subsection{References and status
check}\label{references-and-status-check}

\begin{itemize}
\tightlist
\item
  M. Bogoya, A. Böttcher, M. Ferrari, S. M. Grudsky and S.
  Serra-Capizzano, \emph{Condition numbers of block Toeplitz matrices
  and stability of space-time IgA approximations for the wave and
  Schrödinger equations},
  \href{https://arxiv.org/abs/2608.24151v1}{arXiv:2608.24151v1}, 25
  August 2026. \href{https://arxiv.org/html/2608.24151v1}{Full text}:
  equation (5.22) gives \(B_k\); §5.3, equation (5.32) and the following
  coefficient display give the family and \(C_k\); the paragraph after
  Figure 9 and Remark 4.3 state the unresolved polynomial-growth issue.
\item
  M. Ferrari and S. Gómez, \emph{A matrix-based approach to the
  stability of a space-time isogeometric method for the linear
  Schrödinger equation},
  \href{https://arxiv.org/abs/2506.18859v2}{arXiv:2506.18859v2} (5 May
  2026), abstract and introduction. Its analyzed method uses splines of
  maximal regularity, a different case from cubic \(C^1\) splines.
\end{itemize}

On 2026-09-10, checked the complete August 2026 source, its version
history, and targeted identifier/title/correction searches. Only v1 was
listed, and no later proof or counterexample was found. This short,
bounded follow-up search does not establish exhaustive openness. The
source's more general claim involving only determinant-root counts is
not needed for this particular, explicitly stated family.

The 2026-09-11 submission audit checked five public repositories and all
30 public branch heads, where MF-22 remained Open, and found no prior
full solution in the checked documents or discussions. The
\href{https://github.com/ajt60gaibb/OpenProblemsInNLA/blob/main/references/stepaniants-mf22-2026-09-11/verification/network-check.json}{dated
audit} records the scope and its limits. The source still listed only
v1.
\end{document}
